\section{Information categories and relation to earlier inverses}
\label{sec:comparison}
The stationary results combine a geometric query reduction with an
integral normalization. The conjunction in
Proposition~\ref{prop:rare-query} produces a one-sided boundary test
from raw pooled collision bits, including at ties between compass
directions. The integrated collision strip identifies an obstacle
width sum and hence the ratio of two unknown homothetic footprints.
The area-normalized inverse supplies a second identification argument.
These reductions use the actual collision sensor and the loss specified
in Section~\ref{sec:setting}; the statistical and convex-geometric
principles used after the reductions have established antecedents.

\subsection{Boundary estimation and informative rare responses}
Active boundary estimation is an essential comparison for the finite
theorem, not merely an analogy with passive billiard spectra.
Castro and Nowak~\cite{CN} study adaptive binary classification under
boundary-fragment smoothness and noise assumptions, with matching
orders for their excess-risk criterion. Their input selects a feature
point whose response is a class label. This directly available label
is not our collision bit: a solid start and a free miss have the same
output here. Neither their observation model nor their loss identifies
our finite periodic table. We use an independent physical packing and
binary-tree proof rather than importing a classification lower bound.

Locatelli, Carpentier and Kpotufe~\cite{LCK} give an adaptive strategy
for smooth decision boundaries without prior knowledge of smoothness
and noise parameters. Their membership-query setting and local
line-search/interpolation construction make the closest algorithmic
comparison. In contrast, our smoothness is known, our error is a $C^2$
geometry norm, and the boundary label must first be derived from
pooled collision data. We claim neither a new general
bisection principle nor an advance on adaptation to unknown smoothness.
The reductions in Propositions~\ref{prop:query} and
\ref{prop:rare-query} permit this active-estimation strategy in the
physical collision experiment, at their respective costs.

Yan, Chaudhuri and Javidi~\cite{YCJ} use informative abstention
responses in active learning. Under their assumptions the rarity of
an informative response can yield an inverse-probability cost. Their
observation includes an abstention symbol. Our sensor has only its
original collision bit. The additional argument here is that translated
pooled commands and a conjunction over outward compass candidates
realize the requisite one-sided event for each boundary query. Small
occupation alone would not justify such an improvement: the individual
means in a signed reciprocal difference need not be small. The exterior
zero and interior lower bound are therefore proved for the commands
that are actually sampled.

\subsection{Localization and control precision}

The earlier scalar inverse reconstructed a uniformly accurate occupation
on a full two-dimensional grid. Its constructive upper bound was
$C\nu^{-3}\log(C/(\nu\delta))$ attempts on the same bounded
$C^{6,\beta}$ class, using $\sigma\asymp\nu^{3/2}$ and compensated
support smoothing. That theorem remains in the exact repository archive of A2 v31. The present
procedure uses $O(h^{-1})$ angular lines, a fixed-depth dependency
stencil for each queried point, and the full $C^{6,\beta}$ regularity.
It has a different adaptive command design, not a sharper analysis of
the old uniform design. We have not proved a matching lower bound
for that old design or for uniform spatial preparations.

The resource improvement is accompanied by the explicit localization
and calibration cost in \eqref{eq:precision}. Exact data remain a pair
of functionals on controlled spatial inputs. The lower bound concerns
only their finite binary realizations; it does not assert a low-information
passive invariant. The angle/time/position tolerances shrink with
localization and are not learned from the collision outcomes. Apparatus
motion, precision in representing input locations, and arithmetic
cost are distinct resources. Goldberg and Kwek~\cite{GK} treat
query-coordinate precision as an explicit resource in learning convex
polytopes by exact membership queries. The digital bound in
Theorem~\ref{thm:unk-scale-finite} follows the same accounting principle
for a different class and observation: it includes the additional
integral-measurement centers as well as the fine angular boundary
queries. It does not convert a bit-description bound into a physical
metrology or travel bound.

For completeness the exact v31 manuscript is preserved unchanged
inside the repository archive of the reviewed v33 source. It retains the arbitrary centered-law and general
nontrivial bounded-law inverses, the different-zero-set comparison,
the full fixed-aperture proof, rational interval enclosures, the
uniform-grid construction, and all earlier moment, marked-law and
count-fiber results in their original nested volumes. The more general
randomized law is an exact functional theorem; the finite local query
proved here uses the four-atom compass law. No continuous transition
law is discretized without an error budget. Rational intervals remain
conditional on forcing and survival bounds, not certificates of the
apparatus or physical prior.

The finite global conclusion also retains the known positive patch
margin and the bounded periodic presentation. Exact rational period
relations are distinct from approximate Euclidean generators. Removing
the known margin gives only pointwise eventual recognition, as described
in Section~\ref{sec:periods}. The result does not settle rigidity from
uniformly prepared count germs, exact analytic count germs, marked
length spectra or passive trajectories. These boundaries delimit the
information theorem without changing its topic or weakening any retained
mathematical assertion.

\subsection{Stationary noise versus adversarial implementation}
The new fixed-footprint theorem belongs also to support estimation with
contaminated data. Brunel, Klusowski and Yang~\cite{BKY} observe
continuous random vectors drawn from an unknown convex body and then
contaminated by additive Gaussian noise. Their geometric estimator
avoids Fourier inversion and is analyzed in Hausdorff loss. Here no
random vector or direct inside/outside response is observed: the data
are attempted collision bits at controlled centers. The fixed compact
noise support, its boundary-mass lower bound and the reciprocal reduction
are essential to our polynomial upper bound. The Gaussian observation
and its tail-based analysis are not covered by this theorem, and its
rates are not being compared as rates in a common experiment.

Support addition, inner rolling disks and the Brunn--Minkowski inequality
are classical \cite{Schneider}. Their role here is to remove an unknown
stationary density without learning that density. The positivity set of
the recovered blur, not its detailed values, identifies the footprint
sum. The retained signed high-accuracy query uses the Green bound for exit from a fixed
rectangle; its constants can be large and its finite-depth work grows
logarithmically. The area normalization uses a finite adjoint of the
same killed operator. The direct collision-strip normalization instead
uses Cavalieri's principle and width additivity; only one integrated
pooled forward mean is required. We do not claim these convex-geometric
or random-walk
principles as new abstract theorems.

The two upper bounds cannot be ranked by their attempted-bit exponents
alone. The localized model attains the smaller sample power but requires
shrinking spread and worst-case tolerance. The stationary model pays a
larger, not claimed sharp sample power while keeping its random spread
fixed and allowing the density to be unknown. The footprint and the
relative homothety factor are reconstructed in the strongest theorem.
Exact homothety about the supplied common origin, quantitative
boundary-mass and geometric bounds, per-setting stationarity and
nominal positioning remain apparatus assumptions. No theorem prices their manufacture,
extends the converse to every common noise law, or infers a passive
spectral invariant from these active experiments.
